\documentclass[../../main.tex]{subfiles}
\begin{document}
\section{Excess propagation: audit and a circularity warning}
\label{sec:excess}

This section proves Proposition~\ref{prop:intro-audit} --- the consumption audit showing the
unweighted excess term is inert --- and explains why a direct propagation argument can merely
restate a localized isoperimetric lower bound, motivating the bootstrap mechanism of
Section~\ref{sec:bootstrap}. This is a methodological warning, not a no-go theorem. The static obstruction that
forces the covariance weight (Proposition~\ref{prop:two-tail}) is established in
Section~\ref{sec:qcts}.

\subsection{The consumption audit}

\begin{proposition}[Unconditional integrated excess bound]
\label{prop:trivial-excess}
Let $\mu$ be isotropic log-concave on $\R^n$, $E$ measurable with $\mu(E)=\tfrac12$ and excess
$e_0=e_0(E)$. Then for every $T>0$ and every stopping time $\tau$,
\begin{equation}\label{eq:trivial-excess}
  \E\int_0^{T\wedge\tau}e_t(E)\dd t\ \le\ (1+e_0)\,T .
\end{equation}
\end{proposition}

\begin{proof}
Since $I_{\mu_t}\ge0$, $e_t\le\mu_t^+(E)$ pathwise. By the perimeter supermartingale property
\eqref{eq:perimeter-supermartingale} and Tonelli,
$\E\int_0^{T\wedge\tau}e_t\le\int_0^T\E\mu_t^+(E)\dd t\le T\mu^+(E)=T(I_\mu(\tfrac12)+e_0)$.
Finally $I_\mu(\tfrac12)\le1$: a coordinate halfspace through the median of the first marginal
has perimeter at most $\norm{f_1}_\infty\le1$, the supremum bound for unit-variance log-concave
densities on $\R$ \cite{Bobkov1999LogConcave}.
\end{proof}

\begin{proof}[Proof of Proposition~\ref{prop:intro-audit}]
(a) is Proposition~\ref{prop:trivial-excess}. (b): run the proof of
Theorem~\ref{thm:intro-weighted} in Section~\ref{subsec:consumption} with the weight replaced
by $1$ and the weighted excess propagation input replaced by \eqref{eq:trivial-excess}; every
step goes through with the universal constant $C_0'=2C_0+2C_2(1+e_0)\le2C_0+4C_2$ for
$e_0\le1$, and KLS follows. (c) is Proposition~\ref{prop:two-tail} of Section~\ref{sec:qcts}.
\end{proof}

\begin{remark}[Interpretation of the audit]
\label{rem:audit-interpretation}
Item (b) is not good news about excess propagation; it is a diagnosis of the previously stated
package. An additive excess term whose integral is a priori $O(T)$ is inert next to the $C_0T$
term: it can be deleted without changing the strength of the assumption, so the unweighted
Stein-trace estimate silently carried the entire logical weight of the route. Item (c) shows
this weight cannot be discharged slice-wise. The weighted restatement
(Assumption~\ref{ass:weighted-package}) is the repair: the weighted excess term is no longer
inert, the static obstruction is consistent with it, and the burden is honestly distributed
between a stability estimate (Stein trace) and a propagation estimate (weighted excess), each
now a non-inert stated input of this repaired package. Other proof architectures need not use
this decomposition.
\end{remark}

\subsection{Why direct propagation risks circularity}
\label{subsec:circularity}

\begin{lemma}[Perimeter martingale]
\label{lem:perimeter-martingale}
Within the regularity convention (compact support, smooth data), for a fixed finite-perimeter
set $E$ the process $t\mapsto\mu_t^+(E)$ is a nonnegative martingale, with
\begin{equation}\label{eq:perimeter-sde}
  \dd\,\mu_t^+(E)=\Bigl(\int_{\partial^*E}(x-a_t)\dd\sigma_t(x)\Bigr)\cdot\dd W_t ,
\end{equation}
$\sigma_t$ being the weighted surface measure of $\mu_t$ on $\partial^*E$. The drift of the
perimeter vanishes identically.
\end{lemma}

\begin{proof}
$\mu_t^+(E)=\int_{\partial^*E}F_t(x)\dd\sigma_0(x)$ with $F_t$ the pointwise density process of
\eqref{eq:density-sde}; stochastic Fubini on the compact boundary transfers the driftless SDE
of $F_t$ to the integral.
\end{proof}

\begin{lemma}[Exact excess identity]
\label{lem:excess-identity}
Under compact support, for every $t\ge0$,
\begin{equation}\label{eq:excess-identity}
  \E\,e_t(E)\ =\ e_0(E)+\Bigl[\,I_\mu(p_0)-\E\,I_{\mu_t}(p_t)\,\Bigr].
\end{equation}
\end{lemma}

\begin{proof}
Take expectations in $e_t=\mu_t^+(E)-I_{\mu_t}(p_t)$ and use
$\E\mu_t^+(E)=\mu^+(E)=I_\mu(p_0)+e_0$ from Lemma~\ref{lem:perimeter-martingale}.
\end{proof}

Identity \eqref{eq:excess-identity} is conceptually decisive: excess propagation is not about
the set $E$ at all, except through the random volume $p_t$. It is exactly the question of
lower-bounding the expected isoperimetric profile of the random posterior at the current mass;
the set enters only through the driftless perimeter martingale.

\begin{lemma}[Infima over a fixed competitor family]
\label{lem:inf-martingales}
Let $\mathfrak S$ be any \emph{fixed} family of finite-perimeter sets (countable, or rendered
measurable by the usual selection conventions), and define
$J_t=\inf\{\mu_t^+(S):S\in\mathfrak S\}$. Then $J$ is a supermartingale. Pointwise, on
$\{t<\tau_\eta\}$,
\[
  I_{\mu_t}(p_t)\ge J_t(\eta),
  \qquad
  J_t(\eta):=\inf\bigl\{\mu_t^+(S):
    \mu_t(S)\in[\tfrac12-\eta,\tfrac12+\eta]\bigr\}.
\]
The competitor family defining $J_t(\eta)$ is random and time-dependent, so the first assertion
does \emph{not} imply that $J(\eta)$ or $I_{\mu_t}(p_t)$ is a supermartingale.
\end{lemma}

\begin{proof}
For each fixed $S$, $\mu_t^+(S)$ is a martingale by Lemma~\ref{lem:perimeter-martingale}, so
$\E[J_t\mid\mathcal F_s]\le\mu_s^+(S)$ for every $S$; take the infimum over $S$. The second
pointwise comparison holds because the constrained family at exact mass $p_t$ is contained in
the windowed family when $p_t$ is in the window. No conditional-expectation comparison is
available from this inclusion because the eligible family changes with $t$.
\end{proof}

\begin{remark}[The circularity, formalized]
\label{rem:circularity}
By Lemma~\ref{lem:excess-identity}, excess propagation is equivalent to a lower bound on
$\E I_{\mu_t}(p_t)$. Such a lower bound is itself a Cheeger-type statement for the random
posterior, so an argument that simply inserts a lower bound for the localized profile risks
assuming the phenomenon the program is meant to prove. Lemma~\ref{lem:inf-martingales} supplies
no monotonicity for that profile, because the mass-constrained competitor family changes with
time. Thus ``direct propagation is circular'' should be read as a proof-design warning, not a
formal impossibility result: additional structure could conceivably control the moving
infimum. The external worst-case constant $\hstar_n$ of \eqref{eq:hstar-def} provides one
demonstrably non-circular anchor in a contradiction argument; the next section develops that
bootstrap without claiming it is the only possible anchor.
\end{remark}
\end{document}
